\documentclass[11pt]{article}
\usepackage[a4paper,margin=1in]{geometry}
\usepackage{fontspec}
\usepackage[boldfont=false]{xeCJK}
\setCJKmainfont{FandolSong-Regular.otf}
\usepackage{amsmath}
\usepackage{amssymb}
\usepackage{array}
\usepackage{booktabs}
\usepackage{graphicx}
\usepackage{adjustbox}
\usepackage{enumitem}
\setlist[itemize]{topsep=2pt,partopsep=0pt,parsep=0pt,itemsep=2pt,leftmargin=1.4em}
\setlist[enumerate]{topsep=2pt,partopsep=0pt,parsep=0pt,itemsep=2pt,leftmargin=1.6em}
\usepackage{parskip}
\usepackage{fvextra}
\fvset{breaklines=true,breakanywhere=true,fontsize=\small}
\setlength{\parindent}{0pt}

\begin{document}

\begin{center}
  {\LARGE\bfseries Reaction kinetics}\\[0.35em]
  {\large A Level Chemistry}
\end{center}
\vspace{0.6em}

\section*{Rate of reaction}

The \textbf{rate of reaction} 反应速率 is how fast reactants turn into products. We measure it as the change in concentration (or amount) in each unit of time.

To react, particles must \textbf{collide} 碰撞. The \textbf{collision frequency} 碰撞频率 is how often the particles hit each other. But not every collision leads to a reaction:

\begin{itemize}
\item an \textbf{effective collision} 有效碰撞 has enough energy \emph{and} the correct direction, so a reaction happens.

\item a \textbf{non-effective collision} 无效碰撞 does not have enough energy, or the particles hit at the wrong angle, so nothing happens.

\end{itemize}

So the rate depends on the \textbf{frequency of effective collisions} — how many useful collisions happen each second.

\subsection*{Concentration and pressure}

If you increase the concentration of a solution (or the pressure of a gas), the particles are packed closer together. They collide more often, so there are more effective collisions each second, and the rate goes up.

You can calculate a rate from experimental data — for example, the volume of gas made divided by the time taken.

\section*{Temperature and activation energy}

The \textbf{activation energy} 活化能 ($E_A$) is the minimum energy a collision needs in order to be effective.

The \textbf{Boltzmann distribution} 玻尔兹曼分布 is a graph showing how the energies of the molecules are spread out at one temperature. The curve starts at the origin, rises to a peak, then falls away in a long tail. The total area under the curve is the total number of molecules. Only the molecules to the right of $E_A$ have enough energy to react.

When you raise the temperature:

\begin{itemize}
\item the curve flattens and spreads to the right, so a \textbf{much larger fraction} of molecules now have energy greater than $E_A$.

\item the molecules also move faster and collide more often.

\end{itemize}

The first effect is the bigger one. This is why a small rise in temperature gives a large rise in rate.

\section*{Catalysts}

A \textbf{catalyst} 催化剂 speeds up a reaction but is not used up itself. \textbf{Catalysis} 催化作用 is the name for this action.

A catalyst works by giving the reaction a different \textbf{reaction mechanism} 反应机理 — a new route with a \textbf{lower} activation energy. On the Boltzmann distribution, lowering $E_A$ moves the line to the left, so more molecules now have enough energy. This means more effective collisions each second, and a faster rate.

On a \textbf{reaction pathway diagram} 反应路径图, the catalysed route has a lower energy "hill". The enthalpy change of the reaction, $\Delta H$, is \textbf{not} changed by the catalyst.

There are two types:

\begin{itemize}
\item a \textbf{homogeneous catalyst} 均相催化剂 is in the \textbf{same} physical state as the reactants — for example, an acid catalyst dissolved in a solution of liquids.

\item a \textbf{heterogeneous catalyst} 多相催化剂 is in a \textbf{different} state from the reactants — for example, solid iron speeding up the reaction of gases in the Haber process.

\end{itemize}


\end{document}
